% Mycelium Report Template — OVERVIEW shape
% =============================================================================
% Use when: the audience needs the headline and a sketch of the methods, but
% not the exhaustive table. Target length: 2–5 pages of main text. No
% supplement, no appendix.
%
% Copy to analysis/[name]/reports/[name]-report.tex and replace %%PLACEHOLDER%%
% markers. Numbers and terms must be sourced from .manifest.json; do not
% introduce values that have not been registered there.

\documentclass[11pt,a4paper]{article}

% --- Packages ---
\usepackage[utf8]{inputenc}
\usepackage[T1]{fontenc}
\usepackage{amsmath,amssymb}
\usepackage{graphicx}
\usepackage{booktabs}              % Professional tables
\usepackage{hyperref}
\usepackage[margin=1in]{geometry}
\usepackage{natbib}                % Citations
\usepackage{xcolor}
\usepackage{caption}
\usepackage{subcaption}            % Sub-figures
\usepackage{float}                 % [H] placement

\hypersetup{
    colorlinks=true,
    linkcolor=blue,
    citecolor=blue,
    urlcolor=blue
}

% --- Generated reportable values (see report-values-guide.md / scitexlintr) ---
\IfFileExists{build/report_values.tex}{\input{build/report_values.tex}}{%
  \providecommand{\SciVal}[2]{#1}%
  \providecommand{\SciText}[2]{#1}%
}

% --- Metadata ---
\title{%%TITLE%%}
\author{%%AUTHOR%%}
\date{%%DATE%%}

\begin{document}

% =============================================================================
% TITLE PAGE
% =============================================================================
\maketitle

% =============================================================================
% ABSTRACT
% =============================================================================
\begin{abstract}
%%ABSTRACT%%
% Lead with the question (from the planning brief), the comparison baseline,
% and the primary metric value. State the single biggest caveat. No undefined
% acronyms, no citations, no figure references. 150–250 words.
\end{abstract}

% =============================================================================
% PROBLEM STATEMENT
% =============================================================================
\section{Problem Statement}

%%PROBLEM_STATEMENT%%
% Frame the question in plain English. Name the baseline of comparison.
% State what a good answer would look like. Two short paragraphs is typical
% for the overview shape.

% =============================================================================
% METHODS
% =============================================================================
\section{Methods}

%%METHODS_OVERVIEW%%
% In the overview shape, Methods is a single paragraph: 4–6 sentences that a
% non-specialist could follow. Name the inputs, the analytical operation, the
% comparison baseline, and the primary metric. For load-bearing methodological
% choices, lead with the intuitive explanation before the technical term —
% even in the overview, a sentence of intuition before the formal name pays
% off. No software version numbers, no parameter rationale — those live in
% the analysis directory documentation.

% =============================================================================
% RESULTS
% =============================================================================
\section{Results}

%%RESULTS%%
% Each result is one to two paragraphs: state the question, the finding (with
% the numeric value sourced from the manifest), and the interpretation.
%
% Include exactly one worked example for the first new aggregation analysis
% type the report introduces. Prefer an inline sparkline-style mini-table
% over a full figure. Subsequent claims using the same analysis type reuse
% that example by reference.
%
% Figure template (use sparingly in the overview shape — 1 or 2 figures total):
% \begin{figure}[H]
%     \centering
%     \includegraphics[width=0.75\textwidth]{../outputs/figures/FIGURE_NAME.pdf}
%     \caption{[What it shows]. [How to read it]. [Key takeaway].}
%     \label{fig:LABEL}
% \end{figure}

% =============================================================================
% CONCLUSIONS
% =============================================================================
\section{Conclusions}

%%CONCLUSIONS%%
% Map each conclusion back to the problem statement. State caveats explicitly.
% Distinguish what the data shows from what it suggests. For the overview
% shape, three to five sentences is typical.

% =============================================================================
% NEXT STEPS (optional — delete if not applicable)
% =============================================================================
\section{Next Steps}

%%NEXT_STEPS%%
% Only include if there are clear, actionable follow-up analyses. Delete this
% entire section if there is nothing concrete to say.

% =============================================================================
% PROVENANCE
% =============================================================================
\section{Provenance}

\begin{description}
\item[Analysis directory] \texttt{%%ANALYSIS_DIR%%}
\item[Key scripts] ~
  \begin{itemize}
%%SCRIPT_LIST%%
  \end{itemize}
\item[Git commit] \texttt{%%GIT_COMMIT%%}
\item[Analysis period] %%ANALYSIS_PERIOD%%
\item[Analyst] %%ANALYST%%
\item[Report generated] \today
\item[Manifest] \texttt{.manifest.json}
\item[Compile log] \texttt{.compile-log.md}
\item[Software environment] See \texttt{ENVIRONMENTS\_INSTALLATIONS.md}
\end{description}

% =============================================================================
% REFERENCES (optional — delete if no citations)
% =============================================================================
\bibliographystyle{plainnat}
\bibliography{references}

\end{document}
